% ---------------------------------------------------------------------
% principal/sections/00-introduction.tex
% ---------------------------------------------------------------------

\section{Introduction}
\label{sec:introduction}

State the principal theorem. Then explain how the family is organised:
which part of the argument is self-contained here, and which part is
delegated to a companion. Name the companion explicitly and cite it.

The technical input we need is the estimate of \cite{companion-a},
which we use in the form restated below. A second, independent line of
argument appears in \cite{companion-b}; we do not use it here, but it
gives the same conclusion under weaker hypotheses on the constants.

\begin{theorem}[Principal theorem]
\label{thm:principal}
Let $X$ be a compact metric space and let $\mu$ be a Borel probability
measure on $X$. For every $f \in L^2(\mu)$,
\begin{equation}
\label{eq:principal}
  \int_X f \dd\mu \le \norm{f}_{L^2(\mu)} .
\end{equation}
\end{theorem}

\begin{proof}
Use the companion estimate \cite[Theorem 3.2]{companion-a} to reduce to
the case of bounded $f$, then apply Cauchy--Schwarz.
\end{proof}

The rest of this manuscript proves the reductions that
\cref{thm:principal} needs. \Cref{sec:main-result} carries out the main
argument, and \Cref{sec:conclusion} records consequences.
